
\chapter{Мир}\label{ch:mir}
Внутрироссийские транзакции по~картам \enquote{Мир} остаются
внутри страны, а~криптография опирается на~отечественные алгоритмы.
Из~этих двух условий вырастают маршрутизация через ОПКЦ,
ГОСТ-криптография на~чипе и~собственный TSP в~Mir Pay.

\section{История НСПК}\label{sec:mir-history}

Весной 2014~года санкции против ряда российских банков отрезали
их карты от~процессинга Visa и~Mastercard. Внутрироссийский
карточный оборот оказался в~прямой зависимости от~внешних
операторов~\cite{nspk-about-company,cbr-ps}.

\textbf{АО \enquote{НСПК}} (Национальная система платёжных
карт)~--- оператор платёжной системы
\enquote{Мир}; через её операционную и~платёжную клиринговую систему
проходят внутрироссийские карточные операции всех платёжных систем.
АО \enquote{НСПК} создано 23~июля 2014~года. С~весны 2015~года
внутрироссийские операции по~картам международных платёжных систем
начали проходить через ОПКЦ (Операционный и~платёжный клиринговый
центр; см.~раздел~\ref{sec:mir-opkc}) НСПК: соглашение с~Mastercard
о~переводе внутрироссийского процессинга подписано 30~декабря
2014~года, операционный переход банков завершился к~1~апреля
2015~года; для~Visa полный переход завершён к~лету 2015~года.
Сама платёжная система \enquote{Мир} запущена 15~декабря
2015~года; первые карты представили семь
банков~\cite{nspk-about-company,nspk-mir-10-years}.

С~2018~года карты \enquote{Мир} вышли в~массовый
оборот~\cite{nspk-mir-10-years}. Банк России обязал использовать
их для~ряда бюджетных выплат, а~крупные торгово-сервисные
предприятия~--- принимать~\cite{cbr-nps-faq}.

\highlight{Имя Мир закреплено за~разными сущностями. Мир~--- платёжная система.
НСПК~--- её оператор и~ОПКЦ. Банк России~--- расчётный центр
для~операций между российскими участниками; для~иностранных
участников расчёты ведёт Банк~ВТБ или иные кредитные организации,
привлечённые оператором}~(см.~раздел~1.3.2 Правил
ПС~Мир)~\cite{nspk-about-company,fz-161,nspk-mir-rules-v43}.

\pagebreak[4]
\subsection*{Кобейджинг: Мир--UnionPay и~Мир--JCB}

Кобейджинговая карта связывает внутренний маршрут Мир с~внешним
маршрутом партнёрской системы там, где инфраструктуры Мир нет.
Исторически НСПК выпускала кобейджинговые продукты с~UnionPay и~JCB
(Japan Credit Bureau, японская международная карточная схема);
первая карта Мир--JCB вышла в~августе 2016~года~\cite{unionpayintl-nspk-card-issuance,nspk-jcb-cobadge}.
С~14~марта 2022~года JCB International приостановила обслуживание
российских карт за~рубежом, и~выпуск Мир--JCB российскими банками
фактически прекращён.

Внутри России кобейджинговая карта обрабатывается через сеть Мир,
за~пределами~--- через сеть партнёра, если этот маршрут подключён
конкретным эквайером и~эмитентом. Маршрут на~терминале определяется
не~только страной приёма: приоритет задаёт процедура выбора приложения
(Combination Selection по~EMV Book~1~\cite{emvco-book1}; см.~гл.~\ref{ch:emv})
по~AID (application identifier~--- идентификатор
платёжного приложения на~чипе по~ISO/IEC~7816~\cite{iso-7816-5}),
далее учитываются поддержанные терминалом
ядра\footnote{Ядро (kernel)~--- компонент терминала, реализующий
специфичную для~схемы логику бесконтактного EMV (см.~гл.~\ref{ch:emv}).},
валюта операции и~MCC (merchant category code). Возможен
\emph{третий исход} на~стадии выбора приложения~--- пустой
candidate list: терминал не~поддерживает ни~Mir Kernel, ни~UICS
(UnionPay Integrated Circuit Card Specifications~--- чиповая
спецификация UnionPay), AID-наборы не~пересекаются, и~операция
отклоняется до~GPO либо уходит в~\textit{mag-stripe fallback}
по~конфигурации терминала.
Отказ по~BIN-диапазону (bank identification number, старшие цифры
PAN; ISO/IEC~7812) при~блокировке российских эмитентов
UnionPay International или иностранным эквайером~--- отдельный
шаг: BIN читается из~Track~2 Equivalent Data только после READ
RECORD (нормальный EMV-сценарий) или сразу из~трека при~\textit{mag-stripe fallback};
решение принимается на~эквайере либо в~сети UnionPay,
а~не во~время Combination Selection.
Устройство Мир--UnionPay
разобрано в~разделе~\ref{sec:up-mir}: модели отказа, маршрутизация
по~AID, тестирование комбинации
\enquote{карта × терминал × страна × канал}.

На~2025~год эмиссия Мир--UnionPay ограничена несколькими банками,
а~международный приём резко сузился: 23~февраля 2024~года OFAC внёс
НСПК в~SDN-список, и~зарубежные эквайеры стали отказывать в~обслуживании
карт российских эмитентов независимо от~бренда на~лицевой стороне.

\begin{figure}[t]
  \centering
  \includefiguresvg[width=\linewidth]{assets/figures/ch15-mir/mir-unionpay-routing}
  \caption{Три исхода кобейджа Мир--UnionPay: ОПКЦ~НСПК, UnionPay
  International, пустой candidate list.}\label{fig:mir-unionpay-routing}
\end{figure}
\FloatBarrier

\section{ОПКЦ}\label{sec:mir-opkc}

ОПКЦ (Операционный и~Платёжный Клиринговый Центр)~--- инфраструктурный
узел, через который проходят внутрироссийские карточные транзакции.
С~2015~года и~до~приостановки операций Visa и~Mastercard в~марте
2022~года через ОПКЦ маршрутизировались и~транзакции по~иностранным
карточным схемам: бренд оставался на~лицевой стороне карты,
операционная обработка шла через НСПК~\cite{nspk-about-company,cbr-ps}.

Терминал отправляет авторизационный запрос
в~банк-эквайер, тот направляет его в~ОПКЦ. ОПКЦ идентифицирует
эмитента по~BIN (старшим цифрам PAN, primary account number;
см.~ISO/IEC~7812~\cite{iso-7812-transition}),
маршрутизирует запрос в~банк-эмитент и~возвращает ответ тем же
путём к~терминалу. После авторизации ОПКЦ собирает клиринговую
информацию. В~аббревиатуре ОПКЦ соединены две функции:
\emph{Операционный} центр выполняет онлайн-маршрутизацию
авторизационных сообщений, а~\emph{Платёжный Клиринговый Центр}
формирует клиринговые позиции участников; на~рис.~\ref{fig:opkc-flow}
обе функции совмещены в~одном узле. ОПКЦ НСПК ведёт операционный и~платёжный
клиринг, а~расчёт по~операциям между российскими участниками
завершает Банк России; для~иностранных участников расчётным центром
выступает Банк ВТБ либо иные кредитные организации, привлечённые
оператором (раздел~1.3.2 Правил ПС~Мир)~\cite{fz-161,cbr-732p,nspk-mir-rules-v43}.

\begin{figure}[!htb]
  \centering
  \includefiguresvg[width=\linewidth]{assets/figures/ch15-mir/opkc-flow}
  \caption{Маршрут карточной транзакции через ОПКЦ~НСПК: эквайер,
  обработка, эмитент, расчётный центр.}\label{fig:opkc-flow}
\end{figure}
\FloatBarrier

Протокол обмена между ОПКЦ и~участниками основан на~ISO~8583
(международный стандарт сообщений карточных
авторизаций~\cite{iso-8583-2023}), но НСПК задаёт собственный
профиль сообщений: состав и~интерпретация полей данных определяются
спецификацией НСПК.

Держатель прикладывает карту Мир к~POS-терминалу. Терминал
выбирает платёжное приложение Мир по~AID (раздел~\ref{sec:mir-history})
из~RID (registered application provider identifier, первые
пять байт AID) НСПК
\texttt{A000000658}~\cite{iso-7816-5}, считывает
чиповые данные и~формирует авторизационный запрос.\footnote{Поля
и~значения ниже описывают типичный профиль ISO~8583 в~духе НСПК;
полная спецификация полей и~коды расширений~--- закрытая часть
документации НСПК для~участников.} Эквайер
заполняет сообщение MTI\footnote{MTI~--- message type identifier,
тип сообщения; DE (data element)~--- нумерованное поле данных
ISO~8583~\cite{iso-8583-2023}.}~0100 по~профилю НСПК: DE~2 (PAN,
primary account number, из~диапазона 2200--2204
по~ISO/IEC~7812~\cite{iso-7812-transition}), DE~4 (сумма),
DE~22 (POS Entry Mode~--- способ ввода реквизитов карты
на~терминале), DE~55 (EMV data, включая криптограмму
по~правилам платёжного приложения Мир). Сообщение уходит
в~ОПКЦ по~выделенному каналу. ОПКЦ ищет эмитента по~BIN
в~собственной таблице маршрутизации и~пересылает запрос
на~хост эмитента. HSM (hardware
security module~--- аппаратный модуль безопасности; см.~гл.~\ref{ch:crypto})
эмитента
верифицирует криптограмму по~согласованной с~НСПК схеме
ключей, хост принимает решение и~возвращает MTI~0110
с~кодом DE~39 через ОПКЦ обратно
эквайеру~\cite{nspk-about-company,fz-161}.

Банк, работавший с~Visa или Mastercard до~2022~года, не~мог
переиспользовать хостовый модуль без~адаптации: профили полей, коды
ответов и~правила заполнения отличаются. Хостовая интеграция с~ОПКЦ
требует отдельного проекта подключения с~собственными ключами,
тестовой средой и~сертификацией.

Когда эмитент недоступен, ОПКЦ принимает решение об~авторизации
вместо него по~параметрам, которые эмитент заранее согласовал
с~НСПК. В~Правилах ПС~Мир этот механизм называется
\emph{Сервисом Резервной Авторизации} (раздел~13)~\cite{nspk-mir-rules-v43};
международный аналог~--- \textit{stand-in}. Эмитент задаёт лимит
на~сумму одной операции, накопительные лимиты сумм и~количества
операций по~карте за~заданный период; устаревшие параметры ведут
к~необъяснимым отказам. \highlight{Резервная Авторизация в~ОПКЦ
замкнута внутри российской системы: ни~Visa STIP, ни~Mastercard
\textit{stand-in} processing к~ней отношения не~имеют}. Набор разрешённых
сценариев определяют правила платёжной системы \enquote{Мир}~\cite{nspk-mir-rules,nspk-mir-rules-v43}.

Клиринг в~ОПКЦ работает по~расписанию, зафиксированному в~правилах
платёжной системы~\cite{nspk-mir-rules}. Авторизации, накопленные
за~операционный день, формируют клиринговую позицию, а~расчёт между
участниками завершает Банк России~\cite{fz-161,cbr-732p}. Продавец и~PSP
различают момент авторизации и~момент расчёта. Обещанный срок
зачисления привязан к~моменту попадания транзакции в~соответствующий
клиринговый цикл; одобрение само по~себе срок не~сдвигает.

\section{Mir Pay после 2022~года}\label{sec:mir-pay}

После 2022~года мобильная оплата по~картам \enquote{Мир} держится
на~Mir Pay~--- приложении для~устройств под~управлением Android
с~NFC~\cite{mir-pay-help,nspk-tsp}. Для~держателя это обычная оплата
телефоном; внутри~--- сервис токенизации НСПК, который выпускает
одноразовые платёжные данные в~привязке к~устройству, опираясь
на~криптографию НСПК (вместо VTS или MDES; общая модель сетевой
токенизации~--- в~разделе~\ref{sec:token-providers}).

Собственный TSP (token service provider~--- провайдер сетевой
токенизации; см.~гл.~\ref{ch:tokenization}) у~НСПК появился потому,
что чужой не~подходил. Apple не~открывает Secure Element iPhone
сторонним платёжным системам: Apple Pay работает только с~TSP,
которые Apple интегрировал сам~--- VTS, MDES и~ограниченный набор
локальных провайдеров. После приостановки Apple Pay для~российских
карт в~2022~году канал SE на~iPhone стал для~Мир недоступен.
Использование VTS или MDES вернуло бы внутренний рынок к~зависимости
от~Visa и~Mastercard, ради ухода от~которой НСПК и~создавалась.
Mir Pay работает через HCE (Host Card Emulation) на~Android и~не~требует
доступа к~аппаратному SE~\cite{nspk-tsp,mir-pay-help}.

Mir Pay распространяется как Android-приложение для~смартфонов
с~NFC; минимально поддерживаемая версия Android и~каналы
дистрибуции меняются вслед за~политикой НСПК и~магазинов
приложений. Добавить карту удастся, только если банк-эмитент
подключён к~сервису~\cite{mir-pay-help}. Нативной поддержки Wear OS у~Mir~Pay
на~2025--2026~годы нет: оплата с~часов работает только через
пользовательские обходные сценарии и~схемой не~поддерживается.

Для~банка Mir Pay~--- самостоятельный участок повседневной работы:
интеграция с~сервисом токенизации НСПК, выпуск и~сопровождение
токенов, согласование антифрода и~подготовка службы поддержки
к~отклонённым операциям, проблемам при~добавлении кошелька,
перевыпуску карты, отзыву токена.

\needspace{6\baselineskip}
\section{MirAccept: 3DS-аутентификация в~системе Мир}\label{sec:mir-accept}

MirAccept~--- сервис аутентификации держателя карты для~операций
без~присутствия карты (CNP), который НСПК
эксплуатирует как Directory Server и~инфраструктурную часть
3-D Secure для~системы \enquote{Мир} (см.~гл.~\ref{ch:3ds})~\cite{nspk-miraccept}.\footnote{Актуальная
версия протокола, поддерживаемая MirAccept, относится к~семейству
EMV~3DS~2.x; матрица переноса ответственности
между эмитентом и~эквайером фиксируется правилами платёжной
системы и~сверяется по~официальной документации НСПК.} Архитектурно MirAccept
повторяет общую модель EMV 3DS (см.~раздел~\ref{sec:3ds-three-domains}):
Directory Server поддерживает НСПК, Access Control Server (ACS)~---
банк-эмитент, 3DS Server (компонент EMV~3DS на~стороне продавца
или его PSP, отправляющий запросы аутентификации)~--- продавец или
его PSP (payment service provider~--- провайдер платёжных услуг).

Чтобы получить корректный 3DS-результат по~карте \enquote{Мир},
продавец и~его PSP должны быть зарегистрированы в~MirAccept
и~поддерживать актуальную версию протокола EMV 3DS. Перенос ответственности
\enquote{liability shift} (переход финансовой ответственности
за~спорную CNP-операцию с~эквайера на~эмитента; см.~гл.~\ref{ch:3ds})
при~успешной 3DS-аутентификации работает
по~правилам платёжной системы \enquote{Мир}. Сценарии \textit{frictionless}
(аутентификация без~взаимодействия с~держателем), \textit{challenge}
(с~интерактивной проверкой) и~\textit{attempts response} (формальный
ответ при~невозможности полноценной аутентификации) сохраняются,
но~условия их применения и~матрица ответственности фиксируются НСПК.
Mir Pay и~токенизованные операции
через TSP НСПК встраиваются в~ту же схему 3DS. При~оплате токеном
продавец получает результат аутентификации в~том же формате, что
и~для~нетокенизированной карты.

Подключение MirAccept ломается по~трём шаблонам: 3DS-результат
\enquote{not enrolled} или \enquote{unable to verify} там, где
ожидался успешный \textit{challenge}; отсутствие \textit{liability shift}
при~формально успешной аутентификации; несоответствие версий протокола
между 3DS Server продавца и~ACS эмитента. Все три случая требуют
сверки конфигурации с~НСПК; альтернативного маршрута для~карт
\enquote{Мир} нет.

\section{Тарифы и~interchange системы Мир}\label{sec:mir-interchange}

Тарифную модель платёжной системы \enquote{Мир} устанавливает
оператор~--- НСПК, и~она зафиксирована в~Правилах платёжной
системы и~публикуемых тарифных планах~\cite{nspk-mir-rules,nspk-mir-tariffs}.\footnote{Конкретные
ставки interchange и~комиссии операционных и~клиринговых
услуг НСПК изменяются; для~актуальных значений смотрите
действующую публичную редакцию тарифов.} В~отличие от~Visa Operating Regulations и~Mastercard
TPR (Technical Product Requirements~--- свод
продуктово-технических требований Mastercard), которые задают
глобальные таблицы interchange с~десятками
региональных переменных, тарифы Мир образуют единую
тарифную структуру для~российского рынка.

\textbf{Interchange} (в~Стандарте Системы~---
\emph{межбанковское вознаграждение})~--- плата, которая включается
в~клиринговую позицию участника по~итогам клиринга. Плательщик зависит от~типа операции: при~оплате товаров
и~переводе с~карты на~карту платит эквайрер, при~выдаче наличных,
возврате, запросе баланса и~смене ПИН-кода~--- эмитент (раздел~2.1
Стандарта \enquote{Межбанковские вознаграждения})~\cite{nspk-mir-irf-v211}.
Ставка зависит от~типа карты (Дебетовая, Классическая, Привилегия,
Премиальная, Мир~Суприм и~др.), от~MCC торговой точки и~канала
совершения операции. \textbf{Плата за~операционные, авторизационные,
платежные клиринговые и~прочие услуги}~--- собственная плата
оператора (раздел~2 Тарифов, Приложение~4 к~Правилам), не~относится
к~interchange и~платится участником НСПК как оператору. \textbf{Плата
за~услуги Расчётного центра}~--- отдельная компонента (раздел~3 тех
же Тарифов)~\cite{nspk-mir-tariffs-v414}. \textbf{Торговая уступка},
MSC (merchant service charge),~--- итоговая комиссия эквайера с~продавца;
формируется поверх interchange и~плат НСПК и~согласуется в~договоре
эквайринга, в~Правилах ПС~Мир не~регулируется.

Бюджетные карты \enquote{Мир} (в~Стандарте \enquote{Межбанковские
вознаграждения}, далее~--- Стандарте МБВ~\cite{nspk-mir-irf-v211}, отнесённые
к~продуктовым категориям \emph{Дебетовая} и~\emph{Предоплаченная})
массово выдаются по~обязательной эмиссии и~получают особый тариф:
пониженная ставка межбанковского вознаграждения почти по~всем
категориям ТСП~--- комбинация процентной ставки порядка 0,3--0,5\,\%
от~суммы небольших операций и~фиксированной платы 3--4~руб. за~операции
выше порогового значения. Для~социально значимых MCC (ЖКХ, бюджетные
платежи, медицина, образование) ставка дополнительно ограничена
сверху по~всем продуктам~--- банк обслуживает пенсионеров
и~получателей социальных выплат с~предсказуемой
экономикой~\cite{cbr-nps-faq,nspk-mir-irf-v211}.

\needspace{4\baselineskip}
\subsection*{Социально значимые категории и~тарифные ограничения}

В~период пандемии и~в~последующие годы в~России неоднократно
вводились временные и~постоянные ограничения на~максимальный
размер торговой уступки для~отдельных категорий ТСП~--- так
называемых социально значимых: продукты питания, аптеки,
медицинские услуги, услуги ЖКХ, образование, транспорт.
Регуляторное основание~--- сочетание решений Банка России,
рекомендаций ФАС (Федеральной антимонопольной службы)
и~параметров государственных программ
поддержки малого и~среднего бизнеса; пример~--- программа
возмещения комиссии при~оплате через
СБП (Систему быстрых платежей; см.~гл.~\ref{ch:sbp})~\cite{cbr-sme-sbp-subsidy}.

Для~эквайера это разнесённая тарифная сетка: обычная торговая точка
тарифицируется по~рыночной ставке, точка из~социально значимой
категории~--- по~ограниченной. Для~эмитента~--- уменьшенный interchange
по~соответствующим операциям. Архитектура процессинга держится
на~одном правиле: тариф завязан на~MCC и~дополнительные атрибуты
ТСП в~реестре эквайера, признак самой карты на~ставку не~влияет.
Один и~тот же эмитент в~течение дня получает по~одной и~той же
бюджетной карте разный interchange~--- в~зависимости от~того,
в~каком ТСП прошла покупка.

\needspace{4\baselineskip}
\subsection*{Классификация MCC в~Мир}

Кодификация торговых точек в~системе \enquote{Мир} опирается
на~тот же международный стандарт ISO~18245 (международный реестр
кодов категорий ТСП), что и~в~Visa, Mastercard,
UnionPay~\cite{iso-18245}. Сами коды MCC (merchant category code)~---
четырёхзначные значения от~0742 до~9950, и~различий с~международными
системами на~этом уровне нет: аптеки~--- 5912, продукты~--- 5411,
заправки~--- 5541, авиаперевозки~--- 4511, и~так далее. Различие
живёт в~надстройке: НСПК фиксирует, какие MCC относятся к~социально
значимым категориям для~целей тарифных ограничений, какие MCC
требуют особой обработки антифродом~--- азартные игры, обналичивание,~---
и~какие учитываются в~государственных программах субсидирования
комиссий. При~подключении ТСП в~Мир код MCC напрямую влияет
на~тариф и~на~состав обязательных проверок.

\section{Обязательная эмиссия и~эквайринг}\label{sec:mir-mandatory}

161-ФЗ и~разъяснения Банка России задают две
обязанности: банки выпускают карту \enquote{Мир} для~ряда бюджетных выплат,
а~крупные торгово-сервисные предприятия обязаны её принимать~\cite{fz-161,cbr-nps-faq}.

Если банк обслуживает пенсии, социальные выплаты или зарплату
работников государственных и~муниципальных учреждений, он~обязан
выпускать и~сопровождать карту Мир~\cite{cbr-nps-faq}.

За~выпуском пластика тянется полное сопровождение карты:
мобильный банк, колл-центр, разбор спорных операций, карточные
лимиты, поддержка Mir Pay, маршрутизация внутрироссийского трафика
и~жизненный цикл кобейджингового продукта.

Для~торгово-сервисных предприятий с~ежегодной выручкой свыше
20~млн~рублей приём карт \enquote{Мир} обязателен~\cite{cbr-nps-faq}.
Для~крупной торговой сети, общая выручка которой превышает этот
порог, обязанность распространяется и~на~отдельные точки сети
с~выручкой свыше 5~млн~рублей в~год; санкции за~нарушение включают
административный штраф и~предписание Банка
России~\cite{fz-161}.

Эквайер и~PSP проверяют, что терминальная сеть и~онлайн-канал
не~отфильтровывают BIN-диапазоны карт Мир, что в~TMS (terminal
management system~--- системе централизованного управления
терминальным парком) развёрнуты правильные параметры, что подключение
ТСП учитывает обязательный приём, а~у~службы поддержки в~инструкциях
разведены отказ продавца \enquote{не хотим включать ещё один бренд}
и~нарушение обязательного правила приёма.

На~стороне ТСП это особенно заметно в~больших распределённых сетях.
Требование обязательного приёма доводится до~кассового ПО,
терминального парка, интернет-эквайринга и~до~процедур внутреннего
контроля инцидентов. Иначе одна и~та же компания формально принимает
Мир, но не~проводит операции на~части точек или в~части каналов.

\needspace{6\baselineskip}
\section{Платёжное приложение Мир, SDK, сертификация}\label{sec:mir-certification}

EMV-апплет карты Мир в~методологии НСПК называется
\textbf{MPA (MIR Payment Application)}. Низкоуровневый профиль
MPA на~чипе НСПК публично не~раскрывается: детали доступны
в~документации для~участников.

\highlight{Мир использует собственную мобильную и~терминальную
инфраструктуру. НСПК публикует Mir HCE SDK (software development
kit~--- набор средств разработки) для~токенизации и~платежей по~NFC,
Mir Kernel SDK для~бесконтактного эквайринга на~POS-терминалах
и~MIR CIT Mobile (certification \& integration testing~---
инструмент сертификационного и~интеграционного тестирования)
для~сертификационного тестирования эквайринговой
сети}~\cite{nspk-tsp}.

Криптографические профили MPA, привязка к~алгоритмам
и~сертификационная матрица описаны
в~закрытой документации НСПК. Нормативную базу задают
стандарты ГОСТ~34.10--2018 (электронная подпись),
ГОСТ~34.11--2018 (функция хэширования Стрибог), ГОСТ~34.12--2018
(блочные шифры Кузнечик и~Магма) и~ГОСТ~34.13--2018
(режимы работы блочных шифров)~\cite{gost-34-10,gost-34-11,gost-34-12,gost-34-13};
криптографические алгоритмы подробно разобраны в~главе~\ref{ch:crypto}.

Участие банка или процессинга в~Мир начинается с~допуска по~Правилам
платёжной системы Мир. НСПК публикует сами Правила, форму заявления
на~присоединение и~методику расчёта гарантийного
обеспечения~\cite{nspk-mir-rules}.

Далее банк или процессинг настраивает маршрутизацию, обмен ключами,
тестовые среды и~обработку сообщений по~правилам НСПК. Подключение
к~ОПКЦ, операционные окна и~обработка инцидентов~--- отдельный
план работ.

Доработки в~маршрутизации, терминальном стеке или мобильном сценарии
проверяются в~тестовой среде НСПК: тестовые BIN, тестовые ключи, тестовые
ТСП и~сценарии \textit{stand-in}~--- собственные.

Эти направления двигаются с~разной скоростью и~принадлежат разным
командам. Правила и~договоры ведут бизнес и~комплаенс, хостовую
сертификацию ведёт процессинг, терминальные и~карточные профили~---
эмиссия, вендоры карт и~эквайринговая сеть. Без~единого проектного
управления участие оказывается формально оформленным и~технически
неполным: хостовый канал поднят, а~терминалы или персонализация
чипа всё ещё не~готовы к~промышленной эксплуатации~---\nopagebreak
\highlight{недооценка любого из~направлений упирается в~ключи, тестовые
сертификаты, терминальные ядра, параметры TMS или расписание
обязательного внедрения}.
